\begingroup
\makeatletter\tagsleft@true\let\veqno\@@leqno\makeatother
\newcommand{\IndexSurd}[1]{{}^{#1}\!\surd}
\ArithmeticHead{307}{R}{{\scriptsize ZUR THEORIE DER TRANSSCENDENTEN FUNCTIONEN. I.}}
\NumberedItem{16}
\noindent [S. 51]
\par\indent Für die Dreitheilung hat man
\[
\begin{array}{ccc}
a&b&\dfrac ba=t\\[7pt]
A&B&\dfrac BA=T
\end{array}
\]
\begin{equation}
T^4+(12t-16t^3)T^3+6ttT^3-(16t-12t^3)T+t^4=0\tag*{1)}
\end{equation}
\noindent oder
\[(T-t)^4=16(t-t^3)(T-T^3)\]
\[
\frac Aa=\frac{3t-t^3+T(1-3tt)}{3\{T(1-tt)+t(1-TT)\}}\mathpunct{\raisebox{.5ex}{.}}
\]
\noindent Setzt man
\[T=\operatorname{tang}H^2,\quad t=\operatorname{tang}h^2,\]
\noindent so ist
\[\begin{aligned}
\cos(h+H)&=\IndexSurd4\cos2h.\,\cos2H\\
\sin(h+H)&=\IndexSurd4\sin2h.\,\sin2H.
\end{aligned}\]
\noindent Setzt man
\[T=\operatorname{tang}N,\quad t=\operatorname{tang}n,\]
\noindent so ist
\[\begin{aligned}
3\sin(2N-2n)^2&=4\sin(N+3n).\,\sin(3N+n)\\
\sin(N-n)^4&=\sin4n.\,\sin4N
\end{aligned}\]
\[
\frac{A\cos n}{a\cos N}=\frac{B\sin n}{b\sin N}=\frac{2\sin(N+3n)}{3\sin(2N-2n)}=\frac{\sin(2N-2n)}{2\sin(3N+n)}\mathpunct{\raisebox{.5ex}{.}}
\]
\ClosingRule
\par\medskip\noindent\hfill39\textsuperscript{*}
\clearpage
\endgroup
